\clearpage
\section*{2026-09-15: Parcel measurements need dated boundaries and an explicit land definition}
\textit{Agent-drafted research entry.}

All ten provisional area comparisons from the twelve-parent review match
DOF's final assessment roll for fiscal 2026/27. Printed map dimensions also
reproduce the areas to within 0.1\% for seven of those parents. Purdy,
Weirfield, and Lorimer remain supported by administrative records without an
independent survey calculation. DOF and PLUTO share underlying geography and
area sources, so their agreement alone does not validate the measurements.

Sackett's later parcels can be measured directly. The deed in the March 2026
NYSDEC amendment describes lot 17 as a $265\times100$ rectangle, less a
$10\times10$ notch, plus a $10\times100$ access strip: 27,400 square feet.
Lot 47 is $125\times100=12,500$ square feet. Their combined area is 39,900.
The removed neighboring lot 16 measures 8,600 square feet on the DOF map,
although the FY2027 assessment incorrectly records 12,500. Using the map
dimensions conserves the predecessor's 48,500 square feet; summing the three
assessment entries would add 3,900 square feet incorrectly. A 2022 property
attachment also contains a title survey measuring an earlier lot 17 at
38,499.76 square feet. These dated boundaries alone do not establish the land available at the
original February 2022 decision. The panel's 300-unit measure is an
administrative observation, not proof of an unchanged 2022 proposal.

Noble/Oak's missing lot 10 is explained by a proposed subdivision. The city
notice for the June 9, 2026 Brooklyn CB1 meeting explicitly describes tentative
division of lot 1 into lots 1 and 10, containing the two proposed buildings.
It reports a 173,834-square-foot site. The full DOF tax-lot polygon measures
about 297,134 square feet, near its 299,214-square-foot assessment area;
MapPLUTO's shoreline-clipped polygon measures about 172,691. Spatial metadata
confirms that the supplied PLUTO layer removes water area. Thus the geometry
difference largely reflects a different land definition. The documented
onshore site is the relevant candidate for physical development land; it is
not interchangeable with the full tax-lot area or a calculation of zoning
floor-area entitlements. The child-lot areas require the filed subdivision plan.

\textit{Manual follow-up on Sackett and Noble/Oak.} The live DOB record for
B00665570-I1 reports a 48,500-square-foot zoning lot and 300 units. Its
August 2025 recorded zoning declaration, CRFN 2025000213962, explicitly
includes tentative lots 16, 17, and 47. This explains why the later
39,900-square-foot sale/cleanup boundary need not equal the zoning site.
The live DOB text instead lists 17, 47, and 41; this discrepancy is retained.
A move from the production area's 38,500 to the zoning site's 48,500 would
add 10,000 square feet (26.0\%), but the original decision-date footprint
still needs the dated site plan. DOB lists a ZD1 with a January 20, 2022
creation date, but its contents were not retrievable in this review.
The live I1 record, modified June 2026, is not an immutable 2022 snapshot.

Noble's earlier shoreline-clipped polygon measures about 126,049 square feet,
versus 172,691 in the later PLUTO map. The March 2025 DOF event cites the
July 2024 deed and an October 2024 survey. Manual review of the predecessor
May 2001 deed, reel 5169/page 1555 (Schedule A on image page 3), establishes
boundaries at the centers of Noble, West, and Oak Streets. The 2024 deed,
CRFN 2024000204975, identifies the same premises and includes the seller's
abutting-street rights. Those rights predate 485-x, so the map expansion is
not itself evidence of new post-policy land acquisition. However, the 2024
detailed description starts at West Street's western edge, with east--west
lengths 30 feet shorter than the later map annotations. The cited survey
and filed site plan must reconcile that boundary before the baseline area
is finalized. The 2019 cleanup boundary is a separate two-block site.

The model input should measure the whole parent site's land before its
organization decision, counting that land once. The current pipeline uses
earlier filing-specific PLUTO vintages for historical parents and fixed
PLUTO 23v3\_1 for post-policy parents. Later child lots, cleanup boundaries,
water area, and zoning rights cannot be substituted silently for that input.
The manual review establishes better candidates and the specific remaining
boundary questions; it does not finalize either baseline replacement.

The assessment history makes update lags visible. Chatterton retains 60,979
square feet in FY2026 and changes to 21,133 in FY2027; Fox changes from 2,500
to 10,015 between FY2023 and FY2024. A matching BBL must therefore be paired
with the correct dated boundaries. The frozen public query contains 236
tentative and final records, with 119 final records retained for review. Its
latest underlying extract is May 15, 2026, rather than the September retrieval
date. No production areas, weights, or model fits are replaced in this audit;
the ten provisional comparison means are unchanged.

\small
\textit{Reproduction.} Run \texttt{make dof-subset-review}. The existing
\path{tasks/audits/audit_hdb_mappluto_condo_recovery/} task writes assessment
history, sixteen lot checks, and the waterfront comparison. Its report
\path{report/dof_area_validation.md} supplies sources, formulas, and the
remaining boundary questions. The acquisition audit preserves raw API
responses, query definitions, field documentation, and checksummed downloads
of the additional city/state documents. The source notes in
\path{code/dof_subset_validation_2026-09-15/manual_portal_review.md} under
that acquisition task record exact ACRIS/DOB navigation, documents, pages,
and access limitations. SaveData writes reports as side effects.
\normalsize
